\documentclass{revista}
\usepackage[small, unboxed]{cwpuzzle}

\begin{document}
\Titular%
{TITULO}%
{Santiago González Gómez}%
{pasatempos}%
{SUBTITULO}%

% Descomentar para entrar no modo solución
% No modo solución todos os puzzles saen resoltos
%\PuzzleSolution
\begin{Puzzle}{27}{16}
|[1]P  |{}    |{}    |{}    |{}    |{}    |{}    |{}    |{}    |{}    |{}    |{}    |{}    |{}    |{}    |{}    |{}    |[2]P  |{}    |{}    |{}    |{}    |{}    |{}    |{}    |{}    |{}    |.
|É     |{}    |{}    |{}    |{}    |{}    |{}    |{}    |{}    |{}    |{}    |{}    |{}    |{}    |{}    |{}    |{}    |O     |{}    |{}    |[3]C  |{}    |{}    |{}    |{}    |[4]F  |{}    |.
|N     |{}    |{}    |[5]L  |U     |[6]Z  |{}    |{}    |{}    |{}    |{}    |{}    |{}    |{}    |{}    |{}    |{}    |[7]I  |N     |T     |E     |N     |S     |I     |D     |A     |D     |.
|D     |{}    |{}    |{}    |{}    |E     |{}    |{}    |{}    |{}    |[8]C  |{}    |{}    |{}    |{}    |{}    |{}    |S     |{}    |{}    |R     |{}    |{}    |{}    |{}    |S     |{}    |.
|U     |{}    |{}    |[9]É  |T     |E     |R     |{}    |{}    |[10]G |R     |I     |F     |[11]F |I     |T     |H     |S     |{}    |{}    |N     |{}    |{}    |{}    |{}    |E     |{}    |.
|L     |{}    |{}    |{}    |{}    |M     |{}    |{}    |{}    |{}    |I     |{}    |{}    |U     |{}    |{}    |{}    |O     |{}    |{}    |{}    |{}    |{}    |{}    |{}    |S     |{}    |.
|[12]O |S     |[13]C |I     |L     |A     |D     |O     |R     |E     |S     |{}    |{}    |E     |{}    |{}    |{}    |[14]N |E     |W     |T     |O     |N     |{}    |{}    |{}    |{}    |.
|{}    |{}    |U     |{}    |{}    |N     |{}    |{}    |{}    |{}    |T     |{}    |{}    |R     |{}    |{}    |{}    |{}    |{}    |{}    |{}    |{}    |{}    |{}    |{}    |{}    |{}    |.
|{}    |{}    |A     |{}    |{}    |{}    |{}    |{}    |{}    |{}    |A     |{}    |{}    |Z     |{}    |{}    |{}    |{}    |{}    |{}    |{}    |{}    |{}    |{}    |{}    |{}    |{}    |.
|[15]V |E     |N     |U     |S     |{}    |{}    |{}    |[16]W |O     |L     |F     |R     |A     |[17]M |I     |O     |{}    |{}    |[18]C |{}    |{}    |{}    |{}    |{}    |{}    |{}    |.
|{}    |{}    |T     |{}    |{}    |{}    |{}    |{}    |{}    |{}    |{}    |{}    |{}    |{}    |O     |{}    |{}    |{}    |{}    |A     |{}    |{}    |{}    |{}    |{}    |{}    |{}    |.
|{}    |{}    |O     |{}    |{}    |{}    |{}    |{}    |{}    |{}    |{}    |{}    |{}    |[19]A |M     |A     |{}    |{}    |{}    |U     |{}    |{}    |{}    |{}    |{}    |{}    |{}    |.
|{}    |{}    |{}    |{}    |{}    |{}    |{}    |{}    |{}    |{}    |{}    |{}    |{}    |{}    |E     |{}    |{}    |{}    |{}    |S     |{}    |{}    |{}    |{}    |{}    |{}    |{}    |.
|{}    |{}    |{}    |{}    |{}    |{}    |{}    |{}    |{}    |{}    |{}    |{}    |{}    |[20]I |N     |E     |R     |C     |I     |A     |{}    |{}    |{}    |{}    |{}    |{}    |{}    |.
|{}    |{}    |{}    |{}    |{}    |{}    |{}    |{}    |{}    |{}    |{}    |{}    |{}    |{}    |T     |{}    |{}    |{}    |{}    |{}    |{}    |{}    |{}    |{}    |{}    |{}    |{}    |.
|{}    |{}    |{}    |{}    |{}    |{}    |{}    |{}    |{}    |{}    |{}    |{}    |{}    |{}    |O     |{}    |{}    |{}    |{}    |{}    |{}    |{}    |{}    |{}    |{}    |{}    |{}    |.
\end{Puzzle}

\begin{PuzzleClues}{\textbf{Vertical}}
\Clue{1}{PÉNDULO}{El legado de Foucault en la facultad}
\Clue{2}{POISSON}{Probablemente un pez manchado}
\Clue{3}{CERN}{Centro de investigación europeo más internacional}
\Clue{4}{FASES}{Son lunares o diagramas}
\Clue{6}{ZEEMAN}{Ver doble en cuántica se debe a este efecto}
\Clue{8}{CRISTAL}{Disposición atómica ordenada}
\Clue{11}{FUERZA}{Masa x aceleración}
\Clue{13}{CUANTO}{¿Cuánta energía?}
\Clue{17}{MOMENTO}{Esta revista lo tiene}
\Clue{18}{CAUSA}{Precede al efecto no relativista}
\end{PuzzleClues}

\begin{PuzzleClues}{\textbf{Horizontal}}
\Clue{5}{LUZ}{Onda y corpúsculo}
\Clue{7}{INTENSIDAD}{Voltaje = ... x Resistencia}
\Clue{9}{ÉTER}{Si no se sabe qué es, seguramente sea este}
\Clue{10}{GRIFFITHS}{El gran libro (al menos en EM)}
\Clue{12}{OSCILADORES}{En esencia, todo son ... (pl)}
\Clue{14}{NEWTON}{Físico nacido en Navidad}
\Clue{15}{VENUS}{Afrodisíaco planeta}
\Clue{16}{WOLFRAMIO}{Elemento W}
\Clue{19}{UMA}{Unidad de masa atómica}
\Clue{20}{INERCIA}{Principio de ..., el primero de tres}
\end{PuzzleClues}

Causa • CERN • Cristal • Cuanto • Fases • Fuerza • Griffiths • Inercia • Intensidad • Luz • Momento • Newton • Osciladores • Poisson • Péndulo • UMA • Venus • Wolframio • Zeeman • Éter


\end{document}    
